\section{Related work}
\label{sec:related}

\paragraph{Failure prediction versus intervention value.}
The observation that a failure score is not a decision rule is established, and
we do not claim it.  \citet{interventionparadox2026} show that a critic with
strong offline discrimination can nonetheless degrade agent performance;
\citet{calibrationnotcontrol2026} name \emph{intervention advantage} as the
decision quantity, formalise the mismatch as target error, and introduce
\emph{prefix branching} --- executing candidate actions from an identical
trajectory state --- showing that recalibration improves prediction while
leaving control regret unchanged; \citet{agentforesight2026} supply online
prefix-based warning.  In a non-agentic capacity-constrained placement system,
\citet{fleda2026} likewise observe sub-linear prediction-to-decision transfer.

Our starting point is theirs, and what we add is downstream of it.  We do not
claim the separation between prediction and intervention value; that is
\citet{interventionparadox2026} and \citet{calibrationnotcontrol2026}, and
before them the uplift and welfare-maximization literatures cited in
\cref{sec:main-scope}.  First, the decision object here is the maximum over the
\emph{feasible} set, \eqref{eq:actionable}, under a pre-treatment availability
mask: none of the above defines what happens when the action under study is
unavailable, which is why \cref{thm:refinement} needs the mask as a policy
input.  Second, \cref{thm:ordinal} states the separation as a sharp null rather
than establishing it; the argument is a rearrangement inequality and is claimed
only as the definition of what the deferred policy-ranking study would test,
not as a result.  Third,
prefix branching establishes that intervention advantage is \emph{measurable}
from a reconstructed prefix, and its authors do check the reconstruction:
\citet{calibrationnotcontrol2026} replay each prefix into the environment state
reached by the base trajectory, verify step-by-step agreement, discard any
mismatch (reporting a $100\%$ match rate on retained data), and repeat
branching on an exact-replayable shallow subset (their \S3, B.1, B.7).  What
that check does not do is report fidelity \emph{per state layer}, quantify the
contrast bias that an undetected mismatch would induce, or carry replay
uncertainty into the decision that the measured advantage is used for; and it
does not certify a deployable policy.  Those are the gaps E0-M addresses, and
that \cref{sec:replay-validity} states as an identification limit
with the paired-reference evidence that removes it.  \Cref{thm:safe,prop:certification} address
what it takes to go from a measured advantage to a rule that provably does not
make a running system worse, under simultaneous no-fault, invalid-output and
cost constraints, at a prevalence that differs from the experimental sampling
ratio.

\paragraph{Agent replay, intervention, and recovery.}
\citet{causalagentreplay2026} model a run as a structural causal model and
re-execute forward under a step intervention to attribute failure;
\citet{causalflow2026} pair attribution with counterfactual repair;
\citet{civex2026} map proposed actions to causal queries and return
execute/reject/experiment/abstain verdicts, a verdict space close to our regime
set; \citet{dover2025} debug multi-agent systems by intervening on messages and
plans, \citet{agentcheck2026} supply a reproduce--intervene--mitigate workbench,
and \citet{whentoolsfail2026} perturb tools and measure replanning.  These
establish that branching and repair are implementable.

\paragraph{Harness and runtime optimization.}
The contrast above is about \emph{prediction versus intervention value}; the one
drawn here is about \emph{optimizing the artifact versus deciding when to deploy
it}.  A parallel literature optimises the runtime around a frozen model rather
than the model itself: whole-harness systems search or synthesise harnesses
\citep{metaharness2026,autoharness2026,ahe2026,harnessx2026}, extend the
editable surface to lifecycle interventions \citep{lifeharness2026}, or convert
recurring failures into scoped, regression-checked repairs
\citep{harnessfix2026,selfharness2026}; a learned-editor line post-trains
policies that edit skills, contexts, or harness controllers
\citep{skillr12026,codeskill2026,selfevolve2026,harnessbridge2026,
harnessrlcontrol2026,hase2026}.  \citet{harnessr12026} is the closest of these
to ours: it post-trains a dedicated engineer with online reinforcement learning
to emit lifecycle-wide executable patches --- episode initialisation,
pre-decision context, a pre-action guardrail that may rewrite or veto, and
post-feedback recovery --- rewarded by rerunning the same task batch under the
installed patch.

We claim no priority over executable interventions and do not distinguish
ourselves by completeness: a patch of that kind already carries a host-runtime
contract, can block, rewrite or force an action, can trigger recovery, and is
followed by a frozen target.  The complementary problem is ours.  Given such a
procedure as a candidate regime, \emph{when should it be deployed at a
particular checkpoint}, and can that deployment policy be identified and
certified under availability, clean-run harm, invalid-output and resource
constraints?  The exchange runs one way: a harness engineer supplies elements of
$\Gset$, and \cref{thm:safe} bounds what installing one can cost --- including
the degradation their own matched-baseline comparison records.

One granularity difference is worth stating rather than eliding.  Their overlay
is installed once and governs every subsequent decision in a batch, whereas $g$
is executed at a single registered checkpoint.  What transfers is the hook body
--- the recovery routine, the guardrail --- instantiated as a one-step regime,
not the batch-level overlay.

\citet{gstar2026} is the closest orchestration comparison: it trains an offline
dispatcher to select ready nodes in a multi-agent DAG under resource limits.
Here the object is instead whether to interrupt at one saved scientific-agent
checkpoint, with every feasible registered regime re-executed and its causal
value certified under availability and vector constraints.  G-STAR also notes
that online exploration and training/evaluation cost can be hard to bound.

The two literatures also differ in how an unavailable or broken intervention is
accounted for.  In the harness branch an invalid patch earns reward zero,
conflating ``could not be applied'' with ``was applied and did not help''.  Our
pre-treatment availability mask $A_{ig}$ keeps the two distinct: a masked regime
is never assigned, while a failure caused by an assigned regime remains an
outcome of that regime (\cref{sec:branch}).

\paragraph{Dynamic treatment regimes.}
The formulation of complete regimes and their values is standard in the
potential-outcome literature on dynamic treatment regimes
\citep{murphy2003,murphybingham2009}, and current SMART practice supplies
inference for embedded multistage regimes \citep{tsiatis2024}.  This paper is
the degenerate one-stage case of that machinery.  We claim no methodological
contribution to it; we claim that it is the right frame for the problem and that
the one-stage restriction is what makes the guarantees verifiable rather than
assumed.  Causally, the design is a one-stage dynamic treatment regime;
computationally, policy learning over the feasible set is a full-feedback
constrained contextual-bandit problem, and the two views are complementary.

\paragraph{Automated feature engineering and joint search.}
MORE-FE evolves feature-transformation programs with separate signals for
format, executability, predictive performance, and semantic novelty
\citep{morefe2026}; CoFEH interleaves feature engineering and hyperparameter
optimization to avoid rejecting transformations that look weak under a fixed
downstream configuration \citep{cofeh2026}.  These systems study how candidate
interventions are generated and how search effort is allocated between
interacting AutoML components.  Our question begins after a candidate regime is
frozen: whether it should be deployed at a saved checkpoint, what harm and
vector cost deployment causes, and whether the resulting policy can be
certified.  Neither system is therefore our estimator or safety baseline, and
neither must be a direct comparator to define this adjacent-work boundary.
CoFEH supplies qualitative motivation for F2-FEHPO; MORE-FE's executability and
quality signals motivate keeping unavailability distinct from post-assignment
failure.  Neither validates our checkpoint instrument or identifies our causal
effects.

\paragraph{Partial identification, repetition, and outside evidence.}
Robust best-arm identification under reward contamination shows that
contamination leaves arm means only partially identified.  Decisions then
depend on effective gaps that subtract the contamination allowance
\citep{altschuler2019,mukherjee2021}.  The mean-only condition $\tilde a>2p$ of
\cref{thm:replay-limit} is that phenomenon in our notation.  Treatment choice
under partial identification supplies the decision theory: the conservative
rule used here (certify when the lower end of the identified set clears the
threshold) is standard, and minimax-regret alternatives exist
\citep{manski2000,stoye2012,montielolea2023}.  What neither literature
addresses is repeated sampling at a latent unit.  $R$ restorations of one
checkpoint reveal the depth-$R$ observation moments of that checkpoint's success
probabilities, and the relevant theory is that of truncated moment problems
and mixture identifiability \citep{karlinstudden1966,kemperman1968,schmudgen2017,teicher1963}.
The extremal problem of \cref{thm:finite-repeat} is a worst-case tail mass
over a moment set, the kind of quantity studied for robust risk measures
\citep{rockafellar2000,zhu2009,bertsimas2005}.  When references are available,
combining many replays with a few reference executions is two-phase sampling
\citep{neyman1938,cochran1977}, and prediction-powered inference gives modern
rectified estimators \citep{angelopoulos2023ppi}; we claim nothing new there.
Variance-adaptive confidence bounds for bounded means are established
\citep{maurer2009,waudbysmith2024}; the scaffold uses one as a separate method,
not a contribution.

\paragraph{Credit assignment, grouping, and fixed-state intervention.}
Three families target different objects.  Association-based methods score
steps from observed trajectories or trajectory graphs \citep{g2po2026,salt2026}.
One of them normalises privileged teacher scores within a step and gates them
by their association with return, using prompt groups at completion level and
anchor-state groups at step level \citep{adrs2026}.  Its association statement
concerns the gate's reliability within a batch, not causal identification of a
fixed-state effect; \cref{rem:association} makes the relation to
\cref{thm:ordinal} exact.  AgentOPSD instead scores a turn by how much it
revises a running success belief, read as evidence for eventual success, and
states that the belief is relative support rather than a calibrated
probability; at $\gamma=1$ its recursion is Wald's log-likelihood-ratio
accumulation over turns \citep{agentopsd2026}.  Each estimates from logs a
quantity the instrument executes; the deferred policy-ranking study would carry
an unregistered descriptive secondary ranking checkpoints by a risk-revision score
$r(z_t)-r(z_{t-1})$, the belief-revision analogue (\cref{rem:association}).
Group-based methods compare returns across prompt or
recurring-state groups \citep{gigpo2025}, and history-context inconsistency and
its bias--variance trade-off are studied explicitly \citep{hgpo2026}.  CRAFT
estimates a group-level counterfactual supported by sibling rollouts, and states
that connecting it to a per-trajectory, fixed-state intervention requires an
additional assumption \citep[Sec.~3.1]{craft2026}; that distinction is theirs.
Branching methods fork rollouts at selected steps
\citep{arpo2025,treerl2025,at2po2026}, so alternatives \emph{are} executed
there, but under the training policy rather than as complete oversight regimes
with a frozen continuation and restoration evidence.  A recent self-improvement
loop goes further and scores alternative exploration policies by replaying a
recorded discovery tree, revealing one stored child per node and nothing where
no branch was recorded \citep{dreamrsi2026}: a replay simulator over the recorded
support at depth one, which is the $R=1$ case of \cref{thm:finite-repeat} with an
availability mask induced by the behaviour policy rather than declared before
treatment.  Its online phase resumes saved workspaces without restoration
evidence, so the same-effective-state condition E0 tests applies to it
unchanged.  The adopted estimator
comparison of \cref{app:secondary} contrasts a single fork and context matching
with repeated-branch measurement on the same checkpoints.  It uses generic
diagnostics and is not an implementation of any of these methods.

\paragraph{Separating execution, audit, and consequence.}
Automated research pipelines separate machine-checked proof soundness from the
faithfulness of the formal statement to the intended claim, and from the
research significance of the result.  In CausalSmith the faithfulness audit is
performed by language-model agents, and the significance signal is reported
without being treated as validated \citep[Secs.~5.1--5.3]{causalforge2026}.  The
analogy to our three layers is complementary rather than formal: execution
records which regime ran, audits test a declared scope, and measurements assess
whether the distinction changes decisions.  Where their faithfulness audit is a
language-model judgment whose error rate is not estimated \citep[Sec.~7]{causalforge2026},
the replay audit here is mechanical over declared layers, and its detection and
false-alarm rates on deliberately defective restorations are themselves a
registered measurement (E0-M, \cref{app:e0m-detection}).

\paragraph{Safe policy improvement and pessimism.}
\Cref{thm:safe} is an instance of the pessimistic argument developed in
\citet{laroche2019,kumar2020,xie2021}, and \citet{limasrebro2022} give
lower-confidence-bound policy learning in a one-decision setting.  Dream-RSI
selects the replay-best meta-policy with an in-sample no-worse guarantee on the
fixed history \citep{dreamrsi2026}; \cref{thm:safe} is the held-out, pessimistic
form of that selection step.  What is not
covered there is selection restricted to a \emph{certified-feasible} subset
defined by simultaneously estimated constraints, with a per-context availability
mask, and a guarantee that holds uniformly over prevalence.  That
combination is the content of \cref{thm:safe,prop:certification}, and it is
small.  We present it as such.

\paragraph{Reinforcement-learning readings of the design.}
%% Bibliographic fields for sutton1999, li2006, kearns2002sparse, kakade2002,
%% kearnssingh2002, jin2021, watkins1992, achiam2017, howard1966, pan2022 and
%% gigpo2025 were checked field-by-field against publisher or primary records
%% before these \cite calls were emitted, per README.md "Conventions".
The correspondence in \cref{tab:bellman-dictionary} is a dictionary, not a
method, and each entry has an established home.  The \emph{intervention option}
inside a complete regime --- precondition, action, result-contingent response,
terminal condition --- is an option in the sense of \citet{sutton1999}, with the
availability mask as initiation-set membership and the variable duration of the
option the reason the checkpoint process is semi-Markov until the grid is fixed.
The registered representation $Z=f(H)$ is a state abstraction, and
conditional-mean sufficiency \eqref{eq:conditional-mean-sufficiency} is
$Q$-irrelevance of that abstraction in the taxonomy of \citet{li2006}.
Restore-and-branch is generative-model access --- reset to a state and sample its
successors --- the standard assumption of sample-complexity theory, and the
branch design is the depth-one case of the sparse-sampling planner of
\citet{kearns2002sparse}; what differs here is that the access is tested by E0
rather than assumed, and that it is tested only at the checkpoints the design
samples.  The constrained program of \cref{sec:constraints} is a constrained
decision problem with explicit expected-cost tolerances rather than a
scalarised reward \citep{achiam2017}.  The gated rule of \cref{thm:safe} is
pessimistic value iteration at horizon one on a tabular partition
\citep{jin2021}, in the line cited above \citep{laroche2019,kumar2020,xie2021};
that line supplies conservative value estimation under distribution shift, not
the held-out simultaneous bounds \cref{thm:safe} consumes, which are
\cref{prop:certification}'s union bound and nothing else.  The transfer bound of
\cref{prop:transfer} is the performance-difference lemma of \citet{kakade2002}
combined with the finite-horizon occupancy coupling of \citet{kearnssingh2002}.
Sequential $Q$-learning on this data is named only as what the data cannot
support: its convergence conditions require repeated sampling of every action at
every state \citep{watkins1992}, and depth-one branching never samples a
post-intervention state.  None of these correspondences is claimed as a
contribution, and none of them identifies anything: a Bellman recursion is an
accounting identity for whichever kernel generated the data.  They fix what a
sequential extension would have to add, which \cref{sec:bellman} quantifies, and
that is the only use made of them.

\paragraph{Harm estimands and capacity-constrained oversight.}
\citet{kallus2022} shows that the fraction of units negatively affected is not
point-identified from ordinary experimental marginals.  That result governs
designs assigning one arm per unit; executing every feasible arm at the same
checkpoint identifies instead the fraction of \emph{checkpoints} at which a
regime is harmful in expectation, \eqref{eq:harm-fraction}, which is weaker than
an individual-level claim and stronger than an average.
\citet{oversightcapacity2026} study guard calibration under a bounded human
review capacity; our budget-matched capacity-$q$ comparison shares that
constraint structure, but asks which score should fill the capacity rather than
how large the capacity should be.

\paragraph{Evaluation environments.}
MLGym \citep{mlgym2025}, MLAgentBench \citep{mlagentbench2024}, MLE-bench
\citep{mlebench2024} and Agents' Last Exam \citep{ale2026} are referenced by the
tiers of \cref{tab:tiers}, with MLGym only as historical material in the
current protocol.  None documents restorable mid-trajectory executable state,
which is why \cref{sec:experiments} makes that a gate rather than an assumption.

\paragraph{What we do not claim.}
We do not claim the first failure taxonomy, fault-injection framework,
early-warning monitor, recovery benchmark, or counterfactual re-execution of
agent trajectories, and we do not claim harness editing, program repair, or
workflow synthesis; each exists in the work cited above.  The adjacent
production observations we draw on --- sub-linear transfer in storage
placement, winner's-curse selection feedback in online learning, latency-aware
scheduling, and multi-agent orchestration --- are non-agentic and non-causal:
they motivate the question, and we do not claim that they ground
identification, support a theorem, establish effects at scientific-agent
checkpoints, or justify the one-step restriction.  We do not claim that
prediction and intervention value come apart --- that is
\citet{interventionparadox2026} and \citet{calibrationnotcontrol2026} --- only
that \cref{thm:ordinal} says precisely when they do.  We do not claim a new
safe-reinforcement-learning algorithm, calibration metric, agent architecture,
or released artifact, and we do not claim that any empirical statement here has
been tested: none has.  Nor do we claim the instrument of \cref{sec:evaluation}
exists --- only that \cref{sec:experiments} specifies what would establish that
it does.

\subsection{Coverage of prior branching and replay designs}
\label{app:coverage}

\Cref{tab:coverage} records, for the designs this paper is positioned against,
which elements of the measurement contract each source states or demonstrates,
at the granularity a reader can check.  Every cell is backed by a row of
\texttt{paper/figures/coverage\_sources.csv} naming the version read, the
section, and the supporting excerpt or the bounded search scope; a negative cell
is a statement about the inspected scope only, never about the paper as a
whole.  The row for this paper is a registered plan: no scientific experiment has been
run --- three development exams have --- so nothing in it is demonstrated.

\begin{table}[H]
\centering
\small
\setlength{\tabcolsep}{4pt}
\begin{tabular}{@{}lcccccccc@{}}
\toprule
Design & Mask & Restore & Layers & All arms & Frozen $\rho$ & Harm & Leak bias & Provenance \\
\midrule
ReplayGap \citep{replaygap2026} & -- & D & -- & n.a. & n.a. & -- & -- & P \\
Reset credit \citep{resetcredit2026} & n.a. & n.a. & -- & n.a. & P & -- & -- & ? \\
Shepherd \citep{shepherd2026} & -- & D & -- & -- & ? & ? & -- & P \\
Causal Agent Replay \citep{causalagentreplay2026} & -- & n.a. & -- & n.a. & P & -- & -- & P \\
Calibration Is Not Control \citep{calibrationnotcontrol2026} & -- & D & ? & D & ? & ? & ? & ? \\
Credit-assignment survey \citep{zhang2026ca} & -- & P & P & -- & P & -- & -- & P \\
GiGPO \citep{gigpo2025} & -- & n.a. & n.a. & -- & -- & -- & -- & n.a. \\
HGPO \citep{hgpo2026} & -- & n.a. & n.a. & -- & -- & -- & -- & n.a. \\
ARPO \citep{arpo2025} & -- & ? & -- & -- & -- & -- & -- & -- \\
This paper (registered plan) & P & P & P & P & P & P & P & P \\
\bottomrule
\end{tabular}
\caption{\textbf{Coverage matrix.}  Columns: \emph{Mask} --- which regimes were
applicable is recorded \emph{before} treatment as causal bookkeeping;
\emph{Restore} --- branches resume from restored executable environment state,
not a re-fed text prefix; \emph{Layers} --- restoration fidelity reported per
state layer; \emph{All arms} --- every feasible alternative executed from the
same branch point; \emph{Frozen $\rho$} --- all branches continue under one
fixed policy to a fixed horizon; \emph{Harm} --- the fraction of branch points
harmed in conditional mean is a target; \emph{Leak bias} --- the contrast bias
a failed or partial restoration induces is measured; \emph{Provenance} --- the
reinstated hidden states and seeds are declared (the text re-feed / decoder-state
resume / environment-checkpoint restore distinction of \citealp{zhang2026ca}).
Cells: D~= demonstrated (a measurement is reported); P~= specified or proposed;
--~= not reported in the inspected scope; ?~= scope not readable; n.a.~= not
applicable to that design.  Calibration Is Not Control executes every candidate
\emph{action} from a restored prefix, which is an ``all arms'' D at the level
of primitive actions; our arms are complete oversight regimes followed by
$\rho$.  Sources and scopes: \texttt{coverage\_sources.csv}.}
\label{tab:coverage}
\end{table}

\Cref{tab:abc-mapping} maps this paper against the ABC benchmark-reporting
checklist item by item, with the owner of each outstanding item named.
% Mapping to the Agentic Benchmark Checklist (ABC). Three statuses per row, kept
% apart on purpose: whether the item APPLIES to this design, what IMPLEMENTATION
% evidence exists (an artifact in the tree), and what EMPIRICAL evidence exists
% (none: no scientific run has occurred). A planned or unrun check is never an
% empirical status. Item codes and wording as read from arXiv:2507.02825v5
% (Figures 2--4); the full source-linked ledger with file:line evidence is
% docs/paper_track/W8_AGENTIC_BENCHMARK_CHECKLIST_MAPPING.md.
\begin{table}[p]
\centering\small
\setlength{\tabcolsep}{3pt}
\begin{tabular}{@{}p{0.09\linewidth}p{0.06\linewidth}p{0.55\linewidth}p{0.24\linewidth}@{}}
\toprule
Item & Applies & Implementation evidence (artifact) & Empirical evidence \\
\midrule
\multicolumn{4}{@{}l}{\emph{Task validity (\S4.1, T.1--T.10)}}\\
T.1 & yes & none: versions are digest-pinned, not stated in the prompt; owner: author at R5-2A & none (pre-empirical) \\
T.2 & yes & local inference; bounded identical transport retries (\texttt{retry\_policy.py}) & none \\
T.3 & yes & identity failure ends the run non-qualifying with events preserved (\texttt{e0\_ledger.py}); provider errors recorded, never scored & none \\
T.4 & yes & fresh workspace, session and process per branch; parent truncated before forking; every sibling verified on five layers before treatment (\texttt{e0\_executor.py}) & none \\
T.5 & yes & hidden table bound outside the workspace; candidate scored by the independent restricted-AST adjudicator; oracle outside every workspace; denylisted path segments & none \\
T.6 & yes & hash-bound protocol tree and binding slots; seven measured slots null, no signed \emph{replacement} R5-2A; owner: author & none \\
T.7--T.8 & yes & development pools probed by the ideal solver; final suite not generated; owner: author & none \\
T.9 & yes & registered reference procedure and reachability probe & none \\
T.10 & yes & development-round outlier disclosure; scientific pilot not run & none \\
\midrule
\multicolumn{4}{@{}l}{\emph{Outcome validity (\S4.2, O.a.1--O.i.1)}}\\
O.a.1--O.b.2 & n.a. & outcome is a postcondition on a submitted artifact scored on hidden data, not string matching & --- \\
O.b.3, O.h.2 & yes & difficulty window registered; guessing rate unmeasured; owner: author & none \\
O.c.1--O.c.2 & n.a. & deterministic restricted-AST adjudicator; no LLM judge & --- \\
O.d.1--O.e.3 & n.a. & no unit- or fuzz-test scoring & --- \\
O.f.1 & n.a. & no end-to-end workflow test as outcome & --- \\
O.f.2 & yes & adjudicator deterministic by construction; G4 agreement not run; owner: author & none \\
O.g.1--O.g.3 & n.a. & outcome is a scored artifact, not a target state & --- \\
O.h.1 & yes & frozen tool interface; write-then-submit rule; schema validity reported & none \\
O.i.1 & yes & construct stated as an empirical question; no measurement & none \\
\midrule
\multicolumn{4}{@{}l}{\emph{Benchmark reporting (\S4.3, R.1--R.13)}}\\
R.1 & yes & not released (anonymous submission); owner: author & --- \\
R.2 & yes & harness, restoration surface and digest bindings specified & documentary \\
R.3--R.4 & yes & synthetic tasks in a seed/index namespace ledger; content-identity guard & documentary \\
R.5--R.6 & yes & study table, claim specification, formal estimands & documentary \\
R.7 & yes & failure disposition, forbidden repairs, disclosed limitations & documentary \\
R.8--R.9 & yes & E0-M contrast distortion designed on a disjoint pool (unsigned candidate); not run & none \\
R.10 & yes & registered for E3/E4; descriptive specification for E0-M/E1/E2/E5 adopted & documentary \\
R.11 & yes & status ladders, forbidden promotions, kill rule, both-outcomes reading & documentary \\
R.12 & yes & registered comparator classes (E4 deferred; set unchanged) & none \\
R.13 & yes & no trivial-agent task baseline registered; owner: author & none \\
\bottomrule
\end{tabular}
\caption{Mapping to the Agentic Benchmark Checklist of \citet{abc2025} (43 item codes as read from the v5 figures).  ``Applies'' says whether the item is about something this design does; ``implementation evidence'' names an artifact in the tree or its absence; ``empirical evidence'' is none for every item because no scientific run has occurred.  Documentary evidence is not empirical evidence.}
\label{tab:abc-mapping}
\end{table}

